\documentclass[10pt,a4paper]{amsart}
\usepackage[margin=19mm]{geometry}
\usepackage{amsmath,amssymb,mathtools}
\usepackage[scr=rsfs,cal=euler]{mathalfa}
\usepackage{xfrac}
\usepackage[unicode,hidelinks,bookmarks=false]{hyperref}
\newcommand{\ie}{i.e.\ }
\newcommand{\cf}{cf.\ }
\newcommand{\resp}{respectively\ }
\newcommand{\Subsection}{\subsection}
\providecommand{\nobreakdash}{\nobreak\hbox{-}}
\setlength{\emergencystretch}{2em}
\multlinegap0pt
\usepackage{amssymb}
\usepackage{mathtools}
\usepackage[shortlabels]{enumitem}
\newtheorem{theorem}{Theorem}
\title{A torsion-intersection proof of perfect-cuboid nonexistence on 1072 explicit master-tuple fibers}
\author{Ren\'e Peschmann}
\date{Source: arXiv:2604.28072v1 (CC BY 4.0)}
\begin{document}
\maketitle

\noindent\emph{Source and licence.} A torsion-intersection proof of perfect-cuboid nonexistence on 1072 explicit master-tuple fibers. Version arXiv:2604.28072v1 (CC BY 4.0), distributed by arXiv under the Creative Commons Attribution 4.0 International licence, http://creativecommons.org/licenses/by/4.0/, as recorded in the arXiv licence metadata for this exact version and shown on the abstract page https://arxiv.org/abs/2604.28072v1. Full licence notice: \texttt{licenses/arxiv-2604.28072-CC-BY-4.0.txt}.\par\smallskip

\noindent\textit{Editorial scope.} Counted unit: the article's theorem thm:gpg (source0.tex lines 244--255). The article's complete original proof of it is delivered with it (source0.tex lines 257--280, one \texttt{proof} environment of 24 lines). No mathematics is written, repaired or re-derived: the statement and the proof are the article's own lines, in the article's own notation, and no retained line is substituted or re-flowed.  No further source-local context is needed: every cross-reference of every retained line resolves inside the two delivered spans.  Nothing was omitted from any retained span, so the package carries no omission record.  No retained line contains a citation, so the wrapper carries no bibliography and no bibliography entry.  No retained line contains a figure environment, an image-inclusion command or drawing code, so no image is delivered and the package declares no asset dependency.  Editorial formatting only, in addition: an AMS \texttt{amsart} preamble that restates the article's own declarations and macros used by the retained lines, namely the environments \texttt{theorem} on separate counters, together with the standard packages those lines need.  The retained environments print in this document as Theorem 1 in order of appearance, whereas the article prints the same items with the numbering of its own full document.  The complete native source of the article as distributed in the arXiv e-print for this exact version is bundled unchanged as \texttt{sources/199524/source0.tex}.
\begin{theorem}[$g_+ \cdot g_-$ structure]
\label{thm:gpg}
Let $(a, b, m, n)$ be a master tuple. Set
$$
A = am + bn, \quad B = an + bm, \quad C = am - bn, \quad D = an - bm,
$$
and $g_+ = \gcd(A, B)$, $g_- = \gcd(|C|, |D|)$. Then
$$
g_+ \cdot g_- = \gcd(U_1, U_2) =: g_{\mathrm{scale}}, \quad \gcd(g_+, g_-) = 1.
$$
In particular $g_\pm \mid U_1$, $g_\pm \mid U_2$, and consequently $g_\pm^2 \mid f_1$.
\end{theorem}

\begin{proof}
\emph{Step 1: $g_\pm \mid \gcd(U_1, U_2)$.} The linear combinations
$$
mA - nB = aU_2,\quad nA - mB = -bU_2,\quad
aA - bB = mU_1,\quad bA - aB = -nU_1
$$
show $g_+ \mid aU_2, bU_2, mU_1, nU_1$. Since $\gcd(a, b) = 1$ this gives $g_+ \mid U_2$, and since $\gcd(m, n) = 1$ also $g_+ \mid U_1$. The analogous combinations
$$
mC - nD = aU_2,\quad nC - mD = bU_2,\quad
aC + bD = mU_1,\quad bC + aD = nU_1
$$
yield $g_- \mid \gcd(U_1, U_2)$.

\emph{Step 2: $\gcd(g_+, g_-) = 1$.} Since $a - b$ and $m - n$ are odd, exactly one of each pair $\{am + bn, \, an + bm\}$ and $\{am - bn,\, an - bm\}$ is odd, so both $g_+$ and $g_-$ are odd. Suppose an odd prime $p$ divides both. Then $p \mid A, B, C, D$, hence $p \mid A \pm C = 2am$ and $2bn$, and $p \mid B \pm D = 2an$ and $2bm$. By oddness of $p$, $p$ divides $am, an, bm, bn$. Then $\gcd(m, n) = 1$ forces $p \mid a$ and $p \mid b$, contradicting $\gcd(a, b) = 1$.

\emph{Step 3: $g_+ g_- = \gcd(U_1, U_2)$.} The forward direction follows from Steps 1--2. For the reverse, let $p$ be an odd prime with $p^e \,\|\, \gcd(U_1, U_2)$. Since $\gcd(a, b) = 1$, the entire $p$-part of $U_1 = (a-b)(a+b)$ lies in exactly one factor, so $a \equiv \varepsilon b \pmod{p^e}$ for some $\varepsilon \in \{\pm 1\}$; analogously $m \equiv \delta n \pmod{p^e}$. Reducing the four expressions mod $p^e$:
\begin{itemize}
    \item If $\varepsilon = \delta$: $am - bn \equiv \delta bn - bn = (\varepsilon\delta - 1)bn \equiv 0$ and $an - bm \equiv 0$, so $p^e \mid C, D$, hence $p^e \mid g_-$.
    \item If $\varepsilon = -\delta$: $am + bn \equiv -\delta bn + bn = (\varepsilon\delta + 1)bn \equiv 0$ and $an + bm \equiv 0$, so $p^e \mid A, B$, hence $p^e \mid g_+$.
\end{itemize}
Since $\gcd(U_1, U_2)$ is odd (both factors are odd), the product over all such primes gives $\gcd(U_1, U_2) \mid g_+ g_-$.

\emph{Step 4: $g_\pm^2 \mid f_1$.} From Step 1, $g_+ \mid U_2 \Rightarrow g_+ \mid W_1 U_2$, and $g_+ \mid U_1 \Rightarrow g_+ \mid U_1 V_2$. Therefore $g_+^2 \mid (W_1 U_2)^2 + (U_1 V_2)^2 = f_1$. Same for $g_-$.
\end{proof}

\end{document}
